% !PRINTED 342
\weberpage{342}

One can also derive the designation of the four systems of straight
lines from the single scheme:
\[
\begin{array}{ccc}
(0,\,0) & (1,\,2) & (2,\,1)\\
(1,\,1) & (2,\,0) & (0,\,2)\\
(2,\,2) & (0,\,1) & (1,\,0)
\end{array}
\]
by setting up the rows, the columns, and the combinations of three
that correspond to the positive and negative terms in the formation
of the determinant.

\S.\,92.

\gesperrt{Triple equations.}

The system of the nine points of inflection has the property
that from any two of them one can derive a wholly determinate
third in a rational way. This expresses itself as a property of the
equation of degree $9$ on which the determination of the points of
inflection depends, which has, e.g., the abscissae of these points
as its roots. We set up the following definition:

\gesperrt{An equation without equal roots is called a triple
equation if every two of its roots determine a third,
so that each root of such a triple can be expressed
rationally in terms of the other two, and indeed in such
a way that in an equation}
\weq{1}{x_3 = \Theta\,(x_1, x_2),}
\gesperrt{in which $\Theta$ denotes a fixed rational function, one may put for
$x_1, x_2, x_3$ the roots of any triple whatever
in any order}\footnote{The triple equations of degree $9$ were first treated by
\gesperrt{Hesse} in the memoir: ``Algebraische Auflösung derjenigen Gleichung
$9^{\text{ten}}$ Grades etc.'' in Crelle's Journal, Bd.\,34 (1847).}.

The equation on which the points of inflection depend is a
triple equation of degree $9$. But there are also triple equations
of other degrees, e.g. of degree $7$, to which belong those considered
in the preceding chapter\footnote{Cf. \gesperrt{N\"other}, Mathem. Annalen, Bd.\,15 (1879): ``Ueber die Gleichungen
$8^{\text{ten}}$ Grades und ihr Auftreten in der Theorie der Curven vierter
Ordnung.''}.
